\section{Mechanization}
\label{sec:mechanization}

The artifact, a development in Lean~4 \citep{MouraUllrich21} with Mathlib
\citep{Mathlib20}, contains decision procedures for all three domains that
search for holes in three ways: by bounded enumeration
(\cref{fig:algorithm}), by the summary search of
\cref{sec:complexity-upper}, and, at arity one, by the single-word test of
\cref{app:small-arity}. Their correctness is mechanized end to end for flat
constraints. The development covers arbitrary arity and variance, signed
constraints, satisfiability, the semantic reduction, the alphabet-generic
hole bounds, and both domain transfers. The hole-bound modules include
isolated exchange, induced level maps, normal forms, boundary-pair
propagation, and compression within fixed original dictionaries. The
artifact also contains a certifying solver whose answers are checked by
proved checkers. Lean and Mathlib are both pinned to \texttt{v4.35.0}.

The trusted statement consists of the definitions of trees, the variance
order, flat constraints, satisfaction, and entailment in each domain.
A characterization module connects the path-based order to the greatest
fixed point of the constructor rule and shows that the finite and regular
representations denote the intended tree classes. The kernel checks the
proof terms establishing equivalence with finite search. The procedures
are computable definitions. The proof dependencies and commands for
checking them appear in \cref{app:lean-artifact}.

The final bounds are parameterized by the original reader-state count $N$,
without a preliminary state refinement. Tail and depth use those states
directly; compression uses derivative monitors with $N+j$ states at
level $j$, as accounted for in \eqref{eq:holebound}. The paper and artifact
have the same final bounds. Some intermediate statements and constants
differ; \cref{app:lean} gives the public API, module layout, correspondence,
and these differences. \lean{ElementaryBound.lengthBound} defines
$\holeBound(N,d)$; \lean{ElementaryBound.exists_short_hole} proves the
bounded-depth witness theorem. The back end uses
\lean{RejectedTail.holeLengthBound}, and
\lean{RejectedTail.holeLengthBound_le} proves the closed-form estimate
$H\le2^{2^{1280(N+1)^5}}$.

The module \texttt{Complexity} mechanizes the combinatorial content of
\cref{sec:complexity,app:small-arity}. The summary search is proved sound
and complete, and entailment procedures built on it are proved correct in
each domain. For the label readers,
\lean{Complexity.Summary.card_config_label_le} bounds the number of
configurations so that one configuration fits in
$2^{2^{O(k)}}\cdot\operatorname{poly}(n)$ bits, matching the space bound of
\cref{thm:complexity}; \cref{app:lean-complexity} states the exact bound.
Unlike the nondeterministic search of \cref{prop:summary-search}, which
counts the columns of a derive move to bound its space, the mechanized search
decides reachability of an accepting configuration in a finite graph and
keeps no column counter. The configuration bound supports the space
argument, but the executable search computes the reachable set by
saturation and is not proved to run in that space.
\lean{Complexity.Unary.exists_hole_length_le} proves the bound $N-1$ for every
reader over a one-letter alphabet, and the unary procedures test only the
word up to the first rejected cut. For every arity $n\ge2$ and each domain,
\lean{Complexity.entails_hardnessSystem_iff} and its regular and finite
versions prove the hardness reduction correct, with linearly many literals;
the systems of \cref{prop:cx-length} refute entailment, and
\lean{Complexity.expSystem_shortest_hole} gives their exact shortest hole
length $2(\beta+1)2^\beta$. Not mechanized are the complexity classes
themselves, that is, space-bounded machines, Savitch's theorem, and the
space hierarchy theorem; the PSPACE-hardness of universality; the
polynomial-time tests of \cref{lem:poly-tests}; and the polynomial-time
procedure at arity zero (\cref{prop:arity-zero}).
\Cref{app:lean-complexity} gives the correspondence and the differences.

The artifact is included in
\anon[the supplementary material]{\url{https://github.com/noti0na1/deci-nsse}}.
The NP03 counterexample is derived in full in \cref{app:np03}. A separate
kernel-checked formalization agrees with that argument and is not part
of this artifact.
